The $T_{\text{total}}$ analysis above only accounts for the final
QPU sampling, not the cost of training the angles on quantum
hardware itself. To investigate what would happen if both the angle
training and sampling were performed on actual quantum hardware, we
perform the following analysis.

For each candidate parameter setting,
\begin{equation}
    T_{\text{proxy}} = t_{\text{preprocessing}} + N M \, t_{\text{shot}} + Q \, t_{\text{shot}},
    \label{eq:tproxy}
\end{equation}
where $N$ is the number of COBYLA iterations, the same $K$ from
$T_{AS}$ above, or its cumulative sum across sub-depths for
Interpolation, $M$ is the number of shots used per COBYLA iteration,
$Q$ is the final number of sampling shots once the angles have been
trained and fixed, the same $Q$ already used in $T_{QPU}$ above, and
$t_{\text{shot}}$ is the calibrated per-shot duration defined below.
$t_{\text{preprocessing}}$ is zero for Fixed Angles and Linear Ramp,
the measured transfer-lookup time for Parameter Transfer, and for
Interpolation an added term for its recursive matrix-product-state
evaluations, which the base $NM$ term does not capture.

To determine the recommended strategy, we run these optimizations
on MPS (Aer), each problem instance sampling a random subset of
settings from $N \in \{10,20,40,50,60,75,80,100,150\}$ COBYLA
iterations, $M \in \{10,25,50,100,200,500,1000\}$ shots per
iteration, and $Q \in \{100,200,250,500,1000,2500,5000,10000\}$
final shots for the trained methods, Fixed Angles$^\star$, Linear
Ramp, and Interpolation$^\star$. Fixed Angles$^\dagger$ and
Parameter Transfer require no training and sample only over $Q$, up
to $500\,000$ and $300\,000$ shots respectively. This covers a
total of five strategies.

Running the full $N \times M$ training grid plus the final $Q$
shots on \texttt{ibm\_boston} for every strategy and depth is not
practical, so we simulate that grid with the MPS (Aer) evaluator of
Sec.~\ref{sec:tns} as a hardware proxy and record its cost as
$T_{\text{proxy}}$, then rescale the result to the duration
\texttt{ibm\_boston} would take, once assuming ideal noiseless shots
and once with a noise correction. No hardware samples were collected
for this analysis, only the resource axis is calibrated. We build
these calibrated Window Stickers with the Stochastic Benchmark
framework~\cite{Neira2024}.

We calibrate $t_{\text{shot}}$ two ways, both anchored to the
hardware characterization in Sec.~\ref{sec:hw_training} and
App.~\ref{app:hardware_train}. The noiseless calibration sets
$t_{\text{shot}} = 1/2470\,\text{Hz} \approx 405\,\mu\text{s}$, the
measured \texttt{ibm\_boston} shot rate. The noise-corrected
calibration instead sets $t_{\text{shot}} =
t_{\text{shot}}^{\text{noiseless}} / \sqrt{\gamma}(p)$, where
$\sqrt{\gamma}(p) = F_{CZ}^{\,N_{CZ}/2}$ is the fraction of a raw
shot that constitutes one noiseless-equivalent sample, $N_{CZ} = 2
\times 143 \times p$ is the number of two-qubit gates for depth-$p$
QAOA on the 144-qubit heavy-hex device, and $F_{CZ} = 99.872\%$ is
the median \texttt{ibm\_boston} CZ-gate fidelity. This penalizes
deeper circuits more heavily. Table~\ref{tab:calibration} lists the
resulting per-shot durations and rescaling factors.

\begin{table}[t]
    \centering
    \caption{Calibration factors for the noiseless and
    noise-corrected $T_{\text{proxy}}$ resource axes, by depth $p$,
    measured on \texttt{ibm\_boston}. $s_{\text{nl}}$ and
    $s_{\text{nc}}$ are the factors each method's raw
    $T_{\text{proxy}}$ is rescaled by. Depth $p=1$ has no local
    hardware timing records and is excluded.}
    \label{tab:calibration}
    \begin{tabular}{cccccccc}
        \hline
        $p$ & $t_{\text{hw}}$ ($\mu$s) & $N_{CZ}$ & $\sqrt{\gamma}(p)$ & $t_{\text{nl}}$ ($\mu$s) & $t_{\text{nc}}$ ($\mu$s) & $s_{\text{nl}}$ & $s_{\text{nc}}$ \\
        \hline
        2  & 271.27 & 572  & 0.6933 & 404.86 & 583.97  & 1.4925 & 2.1528 \\
        3  & 271.27 & 858  & 0.5773 & 404.86 & 701.35  & 1.4925 & 2.5855 \\
        4  & 271.27 & 1144 & 0.4806 & 404.86 & 842.33  & 1.4925 & 3.1051 \\
        5  & 280.31 & 1430 & 0.4002 & 404.86 & 1011.64 & 1.4443 & 3.6090 \\
        6  & 275.79 & 1716 & 0.3332 & 404.86 & 1214.98 & 1.4680 & 4.4055 \\
        7  & 298.39 & 2002 & 0.2775 & 404.86 & 1459.19 & 1.3568 & 4.8901 \\
        8  & 280.31 & 2288 & 0.2310 & 404.86 & 1752.49 & 1.4443 & 6.2520 \\
        9  & 271.27 & 2574 & 0.1924 & 404.86 & 2104.75 & 1.4925 & 7.7590 \\
        10 & 280.31 & 2860 & 0.1602 & 404.86 & 2527.81 & 1.4443 & 9.0179 \\
        \hline
    \end{tabular}
\end{table}

This analysis covers the strategies already shown on the Pareto
frontier of Fig.~\ref{fig:recommendation}, Fixed Angles$^\star$,
Fixed Angles$^\dagger$, Parameter Transfer, Linear Ramp, and
Interpolation$^\star$.
We ran the campaign of Sec.~\ref{sec:bench_method} for these
families on 30 heavy-hex-144 MaxCut instances (20 training, 10
test), bootstrap-resampled into stochastic-benchmark
statistics~\cite{Neira2024}, fitting an actionable prescription
curve per method on the training instances and evaluating it on
held-out test instances. At each resource level we take the running
maximum actionable response across families to obtain the Pareto
frontier, crediting a method only where it has direct fitted
evidence rather than flat extrapolation past its last measured
point.

Figure~\ref{fig:pareto-overlay} overlays the resulting actionable
Pareto frontier under both calibrations on a single axis, so the
effect of the hardware noise assumption on the recommended strategy
is visible directly rather than split across separate figures. Line
colour marks which method family holds the frontier at a given
resource level; linestyle distinguishes the noiseless (solid) and
noise-corrected (dashed) calibrations.

At the smallest calibrated budgets, Fixed Angles$^\dagger$, without
any angle reoptimization, holds the frontier under both calibrations,
with occasional Parameter Transfer segments, consistent with these
being among the fastest angle-setting methods in
Fig.~\ref{fig:performance}. As the budget grows, Fixed Angles$^\star$
takes over and holds the frontier across most of the remaining range,
in particular the depth-six variant, which dominates the majority of
the accessible budget on both panels. Interpolation$^\star$ only
claims the frontier in a narrow middle window before Fixed
Angles$^\star$ (p=6) reclaims it at the largest budgets, so its
higher training cost does not translate into a lasting advantage once
that training is itself priced onto the hardware axis, echoing the
diminishing returns from extensive angle optimization already noted
for Fig.~\ref{fig:recommendation}. Linear Ramp briefly holds the
test-panel frontier at intermediate budgets under both calibrations
but does not appear on the training-panel frontier in this sweep.

The noise-corrected calibration pushes the entire frontier to the
right relative to the noiseless one, most visibly at the largest
budgets, since deeper circuits lose a larger fraction of each shot to
gate error and therefore need proportionally more of them to reach
one noiseless-equivalent sample. The ordering of which family holds
the frontier is otherwise preserved between the two calibrations,
shallow unoptimized methods are cheapest at low budget, Fixed
Angles$^\star$ dominates the middle and upper range, and
Interpolation$^\star$ is only transiently competitive. The best
achieved approximation ratio at the largest available budget is
unchanged by the calibration choice ($92.9\%$ train, $91.9\%$ test),
since both calibrations rescale the same underlying best-performing
configuration and only its apparent cost in duration shifts.

\begin{figure}[t]
    \centering
    \includegraphics[width=\linewidth]{figures/pareto_frontier_overlay_hw_calibrations.pdf}
    \caption{Actionable Pareto frontier for 144-node heavy-hex MaxCut
    under the noiseless and noise-corrected hardware calibrations
    (Eq.~\ref{eq:tproxy}), overlaid on a single axis.
    (a) Training instances. (b) Test instances. Line colour indicates
    which method family is the best currently deployable strategy at
    that resource level. Solid lines are the noiseless calibration;
    dashed lines are the noise-corrected calibration.}
    \label{fig:pareto-overlay}
\end{figure}
